% Claims and protocol constants. Every data-derived number in a chapter is \est{claim-id}{field};
% every protocol constant is \proto{key}. Both are defined by files that texgen.py generates into
% ../build/generated/. An undefined id prints a red marker and a LaTeX warning; the claims tests
% turn any unresolved id into a build failure.
\newcommand{\reportmissing}[1]{\textcolor{red}{\textbf{[#1?]}}}
\newcommand{\est}[2]{%
  \ifcsname est@#1@#2\endcsname\csname est@#1@#2\endcsname
  \else\reportmissing{#1.#2}\PackageWarning{report}{Undefined claim #1.#2}\fi}
\newcommand{\proto}[1]{%
  \ifcsname proto@#1\endcsname\csname proto@#1\endcsname
  \else\reportmissing{#1}\PackageWarning{report}{Undefined protocol key #1}\fi}
% Marks an interval computed treating rows as independent (likely too narrow); explained in Chapter 2.
\newcommand{\naivemark}{\textsuperscript{\dag}}
\InputIfFileExists{../build/generated/claims.tex}{}{}
\InputIfFileExists{../build/generated/protocol.tex}{}{}

% Theorem-like environments.
\theoremstyle{definition}
\newtheorem{definition}{Definition}[chapter]
\newtheorem{assumption}{Assumption}[chapter]
\theoremstyle{plain}
\newtheorem{proposition}{Proposition}[chapter]
\newtheorem{lemma}{Lemma}[chapter]
\theoremstyle{remark}
\newtheorem*{remark}{Remark}

% Notation.
\newcommand{\E}{\mathbb{E}}
\newcommand{\Prob}{\mathbb{P}}
\newcommand{\Var}{\operatorname{Var}}
\newcommand{\indep}{\perp\!\!\!\perp}
